El sistema propuesto se sitúa en el cruce de cuatro dominios técnicos: el diseño de circuitos analógicos, los modelos de lenguaje, el aprendizaje por refuerzo y la simulación física de circuitos, articulados sobre una arquitectura de software distribuida. Este capítulo expone los conceptos que sostienen cada uno de esos dominios y, sobre todo, las relaciones entre ellos, porque la aportación del proyecto no reside en ninguno por separado, sino en su integración dentro de un lazo de generación y validación. La exposición sigue el recorrido de los datos en el sistema: desde la descripción en lenguaje natural hasta el netlist validado.



\section{Qué es un agente}

\subsection{Que es un prompt}


\section{Modelos de lenguaje grande}

Un modelo de lenguaje grande es un modelo estadístico, entrenado sobre grandes volúmenes de texto, que estima la probabilidad del siguiente elemento de una secuencia a partir de los anteriores. La generación procede de forma autorregresiva: el modelo produce un elemento, lo incorpora al contexto y repite el proceso, de modo que la salida se construye paso a paso \cite{vaswani2017attention}. Esta capacidad de generación permite que el modelo transforme una descripción en lenguaje natural en una estructura de datos o en código, que es precisamente lo que el sistema necesita en dos puntos: la extracción de los parámetros eléctricos a partir de la descripción del circuito y el cálculo de los valores de los componentes de cada sub-bloque. Internamente, el modelo procesa el texto como una secuencia de unidades llamadas tokens y opera sobre una ventana de contexto de tamaño acotado, que fija cuánta información puede considerar en una sola consulta; esta limitación reaparece más adelante como una de las razones para descomponer el circuito en sub-bloques.

Dos propiedades de estos modelos resultan pertinentes para el diseño. La primera es el aprendizaje en contexto: el modelo ajusta su respuesta a partir de los ejemplos e instrucciones incluidos en la propia consulta, sin necesidad de reentrenarlo \cite{brown2020language}. Esto permite proporcionarle las ecuaciones de diseño de un sub-bloque como contexto y obtener los valores correspondientes. La segunda es la generación de salida estructurada, en ocasiones mediante mecanismos de llamada a funciones, que obliga al modelo a responder en un formato definido ---por ejemplo, un objeto JSON con el tipo de circuito y sus parámetros--- y facilita el procesamiento posterior de la respuesta.

La limitación que motiva el problema de investigación es complementaria a estas capacidades. El modelo genera texto plausible, pero no comprueba por sí mismo lo que produce \cite{ji2023survey}. En el dominio de los circuitos, esto se traduce en componentes inexistentes, restricciones físicas violadas y salidas eléctricamente inconsistentes que rara vez funcionan al primer intento \cite{analogcoder,analogagent}. El modelo no dispone de un mecanismo interno que contraste su propuesta con el comportamiento físico del circuito; produce una respuesta que parece correcta según los patrones del lenguaje, no según las leyes de la electrónica. Esta brecha entre plausibilidad lingüística y corrección física es la que justifica incorporar una verificación externa al modelo, y es el punto del que parte el diseño del sistema: la simulación aporta la comprobación que el modelo no realiza.

La generación del modelo no es determinista. En cada paso, el siguiente elemento se elige de una distribución de probabilidad, y un parámetro de temperatura regula cuánta variación se admite en esa elección. Una temperatura baja produce salidas más estables y, al fijarla, hace repetible la generación, condición necesaria para comparar ejecuciones durante la evaluación. Esta misma variabilidad sustenta la estrategia de referencia frente a la que se contrasta el sistema. El enfoque \textit{best-of-N} consiste en pedir al modelo $N$ propuestas independientes para la misma descripción, evaluarlas y conservar la mejor \cite{chen2021evaluating}. Es una forma directa de aprovechar el muestreo, pero no incorpora realimentación: cada propuesta se genera sin información sobre los fallos de las anteriores. El sistema propuesto se mide contra esta línea base a igual número de llamadas al modelo, lo que permite distinguir la mejora atribuible al lazo de ajuste de la que se obtendría por el solo hecho de muestrear más.

\section{LangGraph}

\section{LangFuse}

\section{pyspice}



\section{Sistemas multiagente y orquestación}

Hay que modificar el Map reduce con el Master-Workers.
primero hablarlo 

\section{Arquitectura de software: principios y patrones de coordinación}

La forma del sistema se apoya en principios reconocidos de la ingeniería de software.

\subsection{Separacion de intereses}
cada parte del sistema atiende una preocupación distinta, de modo que las responsabilidades no se mezclan y cada una puede evolucionar por separado \cite{pressman2010software}. El proyecto aplica este principio al dividir el sistema en un subsistema cliente, que gestiona el acceso de los usuarios, y un subsistema servidor, que aloja la lógica de diseño de circuitos. La gestión de usuarios y la generación de circuitos son preocupaciones independientes y se despliegan por separado.

\subsection{Ausencia de estado en el componente a escalar} Una instancia sin estado no guarda información entre solicitudes, de modo que cualquier instancia puede atender cualquier solicitud y el sistema admite replicar el servidor detrás de un balanceador de cargas para repartir la demanda \cite{fielding2000architectural}. Este escalado horizontal es la razón por la que el subsistema servidor se diseña sin estado.
\subsection{La ausencia de estado entre solicitudes convive con una exigencia opuesta dentro de cada solicitud}
El lazo de ajuste necesita conservar el estado entre iteraciones, porque cada ciclo parte del resultado del anterior. El sistema resuelve esta tensión mediante un mecanismo de \textit{checkpointer} que persiste el estado de la generación en una base de datos tras cada paso, de modo que el estado de una solicitud en proceso queda disponible para todos los agentes a lo largo de las iteraciones sin residir en la memoria de una instancia concreta. La coordinación entre agentes se modela como un grafo de estado, en el que los nodos son los agentes y las aristas el flujo de control, y sobre el que se apoya la persistencia por \textit{checkpointer} \cite{langgraph}. Estos conceptos ---separación de intereses, ausencia de estado, persistencia explícita y grafo de estado--- son los que el capítulo de diseño aplica para justificar la estructura concreta del sistema.

\section{El diseño de circuitos analógicos como proceso iterativo}

El diseño de un circuito analógico no se resuelve con un procedimiento cerrado que, a partir de una especificación, devuelva una única respuesta correcta. El proyectista parte de un conjunto de metas eléctricas ---una ganancia, una frecuencia de corte, un punto de polarización--- y propone una topología y unos valores de componentes que, según las ecuaciones de diseño, deberían cumplirlas. La propuesta se verifica mediante simulación o medición, se comparan los resultados con las metas y, cuando hay desviación, se ajustan los valores y se repite el proceso. Este ciclo de proponer, simular, revisar y ajustar se recorre varias veces antes de que el circuito cumpla las metas \cite{sedra2015microelectronic}.

Las ecuaciones de diseño son el vínculo entre el dominio de la especificación y el de la implementación. Para un amplificador no inversor con amplificador operacional, por ejemplo, la ganancia se relaciona con el cociente de dos resistencias; para un filtro RC de primer orden, la frecuencia de corte queda determinada por el producto de una resistencia y una capacitancia. Estas relaciones permiten calcular un valor inicial, pero no garantizan que el circuito real cumpla la meta, porque el modelo ideal ignora efectos que la simulación sí captura: la carga entre etapas, las no idealidades de los componentes y la interacción entre sub-bloques. De ahí que la simulación funcione como árbitro: es el mecanismo que comprueba si los valores calculados producen el comportamiento esperado.

El proyecto se restringe a circuitos analógicos de corriente alterna y corriente directa construidos con componentes comerciales ---resistores, capacitores, inductores, fuentes, diodos, transistores y amplificadores operacionales descritos por macromodelo---, en contraste con la línea de trabajo que diseña a nivel de transistor para un proceso de fabricación específico \cite{skywater_pdk}. Esta distinción no es accesoria. El diseño a nivel de componente comercial trabaja con dispositivos cuyos modelos ya existen y cuyos valores se eligen de catálogos discretos, mientras que el diseño a nivel de transistor dimensiona dispositivos sobre un nodo tecnológico. La frontera entre ambos niveles delimita el alcance del proyecto y, como se argumenta en el estado del arte, marca el espacio que la propuesta ocupa.


\section{Síntesis del marco teórico}

Los conceptos expuestos forman una cadena que recorre el sistema de extremo a extremo. El dominio del diseño analógico fija el problema y explica por qué la simulación debe actuar como árbitro. El modelo de lenguaje aporta la capacidad de pasar del lenguaje natural a una propuesta de circuito, y su incapacidad de autoverificarse motiva el resto del diseño. La organización multiagente reparte ese trabajo en responsabilidades acotadas y, frente a la generación repetida sin realimentación, sustituye el muestreo ciego por un ciclo informado por la simulación. El aprendizaje por refuerzo da marco formal a la decisión que convierte las métricas de la simulación en una acción sobre el circuito. La simulación SPICE produce esas métricas, que son la señal objetiva que cierra el lazo. Y los principios de arquitectura sostienen el conjunto como un sistema que separa preocupaciones, escala y conserva el estado donde debe conservarlo. La articulación de estos elementos dentro de un único lazo de generación y validación, a nivel de componentes comerciales, es lo que el capítulo de diseño desarrolla.

